The 1975--76 material comprises 34 catalogue records and 56 scanned pages,
\GSpage{P000603}--\GSpage{P000658}. The occurrence index below contains
@@ITEMS@@ distinct items. Record 378 is a year heading and has no invented
mathematical content. Catalogue scan envelopes are retained for talks;
the narrower page references on items identify the actual passages.
This distinction matters because several scans contain successive
speakers, and the catalogue does not consistently identify interventions.
In particular, record 383 is catalogued under Novikov, but its manuscript
heading names Drinfeld. Shnirelman's substantial spectral discussion is
embedded within record 391, catalogued under Krichever.

\begin{center}
\begin{tabular}{lrrrrrrr}
\toprule
Class&A&B&C&D&E&F&Total\\
\midrule
Items&@@A@@&@@B@@&@@C@@&@@D@@&@@E@@&@@F@@&@@ITEMS@@\\
\bottomrule
\end{tabular}
\end{center}

The principal mathematical line is the passage from stationary integrable
equations to commuting ordinary differential operators, spectral curves,
and line bundles. The autumn talks explain finite-gap and rank-one
constructions; the February--March discussions articulate the role of
filtered modules, the punctured-curve ring, acyclic line bundles, and
translation on the Jacobian. Drinfeld's 1977 note is an unusually precise
publication match, while Krichever's subsequent treatment separates
higher rank from merely choosing two operators of noncoprime orders
\citep{drinfeld1977commutative,krichever1978commutative}.
Earlier finite-gap occurrences remain essential context, but a recurrence
of a subject is not automatically the recurrence of an identical problem.

Polyakov's sequence moves from perturbative quantum fields and particle
models to gauge theory, static charges, canonical reduction, and a
lattice integral. The corrections repeatedly concern the distinction
between a group and its Lie algebra, gauge-invariant quantities and
gauge-dependent potentials, formal functional integrals and finite
measures, and strong-coupling lattice results and the four-dimensional
continuum confinement problem. The lattice construction was already
published before these talks \citep{wilson1974confinement}; that historical
fact is separated from the unresolved continuum claim. The spring
Hamiltonian calculation does not derive confinement simply by averaging
an instantaneous Coulomb kernel.

The other lines are not marginal appendices. Manin's arithmetic discussion
contains exact formulas whose numerical normalizations can be tested;
Kirillov's talks connect local Lie brackets, deformation, and
supergeometry; Roiter's arithmetic root criterion has an exact later
match in Proposition 4 of his 1978 paper \citep{roiter1978roots}.
Novikov's radial inverse problem is identifiable as Brodskii's scalar
model, rather than a theorem recovering every parameter of an arbitrary
elastic solid. The concluding algebraic discussion belongs to the
period of the solution of Serre's projective-module problem.

Two bounded developments are supplied. For the Fuchs--Perchik question,
the relative Hamiltonian Euler product is normalized, its finite cutoff
is proved, and two exact algorithms agree through weight 300. There are
135 nonzero positive-weight Euler coefficients in that range, and the
notebook's lower bound of 57 classes is recovered. These finite facts do
not prove infinite nonvanishing. The relevant modern computation retains
the exact infinitude question \citep[\S2]{mohammadzadehfuchs2014cohomology}.
For Shnirelman's Abel spectral-phase expression, a Fourier convergence
criterion, a complete circle calculation, and a qualified sphere
counterexample identify the role of a uniform bound. Neither development
is presented as a solution of its unbounded general question or as a
priority claim for the elementary arguments.

There are @@PROBLEMS@@ indexed problems and @@CORRECTIONS@@ correction or
qualification entries. Class F marks an identification limit, not an
undeveloped research programme; no class D is assigned merely because a
search failed. Literal errors, omitted hypotheses, and publication matches
are assessed independently. The report and source ledgers explain the
remaining uncertainty, distinguish manuscript submission from public
circulation, and record which references were read in full and which
were checked only bibliographically. All present-status statements are
bounded by the documented searches as of 10 November 2026.
